\documentclass{article}
\usepackage{graphicx}
\usepackage{amsfonts}
\usepackage{amsmath}
\usepackage{amsthm}
\usepackage{tikz}

\theoremstyle{definition}
\newtheorem{theorem}{Theorem}[section]
\newtheorem{corollary}[theorem]{Corollary}
\newtheorem{proposition}[theorem]{Proposition}
\newtheorem{lemma}[theorem]{Lemma}
\newtheorem{definition}[theorem]{Definition}

\theoremstyle{remark}
\newtheorem{remark}[theorem]{Remark}
\newtheorem{example}[theorem]{Example}
\newtheorem{claim}[theorem]{Claim}

\title{Divide and Conquer}
\author{Sarah Liu}
\date{January 2026}

\begin{document}

\maketitle

\section{Master Theorem}
\begin{theorem}[Master Theorem]
    Let $a\geq 1$ and $b>1$ be constants, $f(n)$ be a function, and $T(n)$ be defined on nonnegative integers by the recurrence $T(n)\leq a\cdot T(\frac{n}{b})+f(n)$, where $n/b$ can be $\lceil \frac{n}{b}\rceil$. Let $d=\log_b a$. Then:
    \begin{enumerate}
        \item If $f(n)=O(n^{d-\epsilon})$ for some constant $\epsilon>0$, then $T(n)=O(n^d)$.
        \item If $f(n)=O(n^{d}\log^k n)$ for some constant $k\geq0$, then $T(n)=O(n^{d}\log^{k+1} n)$.
        \item If $f(n)=O(n^{d+\epsilon})$ for some constant $\epsilon>0$, then $T(n)=O(f(n))$.
    \end{enumerate}
\end{theorem}

\end{document}
